% ID: SEL-D0025-OA-E03
% ESTADO: cotejada. Segunda pasada visual completa y borrador compilado revisados.
% Fuente: Selectividad/matematicas_aplicadas/suplemente_junio_EXAMEN 4-ANDALUCÍA.pdf
% SHA-256: c1a3a7a37d38b04359115c64762b1f8a543e4d66155338931de84ef6c304da2c
% Página: 1; opción A; ejercicio original 3.

En un departamento de una Universidad hay 8 profesores y 14 profesoras. Se quiere constituir una comisión formada por 2 miembros del departamento, elegidos al azar.
\begin{enumerate}[label=\alph*)]
\item \textbf{(0.75 puntos)} \textquestiondown{}Cuál es la probabilidad de que sean profesoras?
\item \textbf{(1 punto)} Calcule la probabilidad de que la comisión esté constituida por un profesor y una profesora.
\item \textbf{(0.75 puntos)} Halle la probabilidad de que en la comisión no haya ninguna profesora.
\end{enumerate}
